% Options for packages loaded elsewhere
\PassOptionsToPackage{unicode}{hyperref}
\PassOptionsToPackage{hyphens}{url}
\PassOptionsToPackage{dvipsnames,svgnames,x11names}{xcolor}
\documentclass[
  11pt,
  a4paper,
]{article}
\usepackage{xcolor}
\usepackage[margin=2.6cm]{geometry}
\usepackage{amsmath,amssymb}
\setcounter{secnumdepth}{5}
\usepackage{iftex}
\ifPDFTeX
  \usepackage[T1]{fontenc}
  \usepackage[utf8]{inputenc}
  \usepackage{textcomp} % provide euro and other symbols
\else % if luatex or xetex
  \usepackage{unicode-math} % this also loads fontspec
  \defaultfontfeatures{Scale=MatchLowercase}
  \defaultfontfeatures[\rmfamily]{Ligatures=TeX,Scale=1}
\fi
\usepackage{lmodern}
\ifPDFTeX\else
  % xetex/luatex font selection
  \setmainfont[]{Libertinus Serif}
  \setmonofont[Scale=0.85]{Latin Modern Mono}
  \setmathfont[]{Latin Modern Math}
\fi
% Use upquote if available, for straight quotes in verbatim environments
\IfFileExists{upquote.sty}{\usepackage{upquote}}{}
\IfFileExists{microtype.sty}{% use microtype if available
  \usepackage[]{microtype}
  \UseMicrotypeSet[protrusion]{basicmath} % disable protrusion for tt fonts
}{}
\makeatletter
\@ifundefined{KOMAClassName}{% if non-KOMA class
  \IfFileExists{parskip.sty}{%
    \usepackage{parskip}
  }{% else
    \setlength{\parindent}{0pt}
    \setlength{\parskip}{6pt plus 2pt minus 1pt}}
}{% if KOMA class
  \KOMAoptions{parskip=half}}
\makeatother
\usepackage{longtable,booktabs,array}
\newcounter{none} % for unnumbered tables
\usepackage{calc} % for calculating minipage widths
% Correct order of tables after \paragraph or \subparagraph
\usepackage{etoolbox}
\makeatletter
\patchcmd\longtable{\par}{\if@noskipsec\mbox{}\fi\par}{}{}
\makeatother
% Allow footnotes in longtable head/foot
\IfFileExists{footnotehyper.sty}{\usepackage{footnotehyper}}{\usepackage{footnote}}
\makesavenoteenv{longtable}
\setlength{\emergencystretch}{3em} % prevent overfull lines
\providecommand{\tightlist}{%
  \setlength{\itemsep}{0pt}\setlength{\parskip}{0pt}}

\providecommand{\E}{\mathbb{E}}
\providecommand{\N}{\mathbb{N}}
\providecommand{\Z}{\mathbb{Z}}
\providecommand{\R}{\mathbb{R}}
\providecommand{\eps}{\varepsilon}
\providecommand{\ceil}[1]{\left\lceil #1\right\rceil}
\providecommand{\floor}[1]{\left\lfloor #1\right\rfloor}
\AtBeginDocument{\let\setminus\smallsetminus}
\usepackage[most]{tcolorbox}
\newtcolorbox{warning}{colback=red!5,colframe=red!65!black,boxrule=0.8pt,arc=2pt,left=6pt,right=6pt,top=5pt,bottom=5pt}
\usepackage{etoolbox}
\AtBeginEnvironment{longtable}{\small}
\setlength{\emergencystretch}{3em}
\usepackage{fancyhdr}
\pagestyle{fancy}
\fancyhf{}
\fancyhead[L]{\small\itshape Candidate result 2510.11311\_\_04 --- GPT-6 Astra, awaiting review}
\fancyhead[R]{\small\thepage}
\renewcommand{\headrulewidth}{0.2pt}
\usepackage{bookmark}
\IfFileExists{xurl.sty}{\usepackage{xurl}}{} % add URL line breaks if available
\urlstyle{same}
\hypersetup{
  pdftitle={A candidate counterexample to Question 6.4 in Extending Thomassen's conjecture to directed graphs},
  pdfauthor={github.com/graph-theory-AI --- machine-generated candidate writeup},
  colorlinks=true,
  linkcolor={blue!50!black},
  filecolor={Maroon},
  citecolor={Blue},
  urlcolor={blue!50!black},
  pdfcreator={LaTeX via pandoc}}

\title{A candidate counterexample to Question 6.4 in Extending
Thomassen's conjecture to directed graphs}
\author{\href{https://github.com/graph-theory-AI}{github.com/graph-theory-AI}
--- machine-generated candidate writeup}
\date{17 September 2026}

\begin{document}
\maketitle
\begin{abstract}
The anti-directed C4 is not Eulerian-avoidable: Eulerian oriented graphs
can have arbitrarily large minimum out-degree while having no
anti-directed-C4-free subdigraph of minimum out-degree 2. This candidate
disproof was generated by GPT-6 Astra in a single model pass for the
Graph Theory AI project. It was selected from the campaign because the
model marked it \texttt{would\_publish:\ true}; it has not been checked
by the adversarial referee, a human mathematician, or a novelty search.
\end{abstract}

\section{The problem and the claimed
result}\label{the-problem-and-the-claimed-result}

This document concerns
\href{https://graph-theory-ai.github.io/graph-conjectures/arxiv/2510.11311__04/}{Question
6.4} from \emph{Extending Thomassen's conjecture to directed graphs}.

\subsection{Problem statement}\label{problem-statement}

Question. Is every orientation of \(C_{4}\) Eulerian-avoidable?

\subsection{Catalog context}\label{catalog-context}

Restricting to Eulerian digraphs (where each vertex has equal out- and
in-degree) may make avoidability questions more tractable. The authors
ask whether the specific case of \(C_4\), which is open in the general
setting, becomes resolvable in the Eulerian host digraph setting.

\subsection{Claimed result}\label{claimed-result}

\textbf{Model verdict: disproved} (confidence: high).

The anti-directed C4 is not Eulerian-avoidable: Eulerian oriented graphs
can have arbitrarily large minimum out-degree while having no
anti-directed-C4-free subdigraph of minimum out-degree 2.

\emph{The model's own caveat:} The counterexamples have very large
degree spread; this does not classify the other orientations of C4.

\section{Model writeup}\label{model-writeup}

The text below is the model output verbatim, apart from moving its
machine-readable verdict block into the summary above and shifting
Markdown heading levels for this document. Mechanical TeX defects and
equation tags were normalized where needed for compilation.

\subsection{A counterexample for the anti-directed
orientation}\label{a-counterexample-for-the-anti-directed-orientation}

Let \(F=\vec K_{2,2}\), with arcs \[
u\to x,\qquad u\to y,\qquad v\to x,\qquad v\to y.
\] All digraphs below are finite and simple, and copies need not be
induced.

\textbf{Theorem.} For every integer \(d\ge 2\), there is a strongly
connected Eulerian oriented graph \(D_d\) with \[
\delta^+(D_d)=d
\] such that no nonempty subdigraph of \(D_d\) has both minimum
out-degree at least \(2\) and no copy of \(F\).

Thus the answer to the question is \textbf{no}.

The construction combines a small terminal bottleneck with a long
sequence of regular, high-girth bipartite layers. Avoiding \(F\) bounds
the number of surviving vertices immediately before the bottleneck. The
high-girth layers then force these numbers to decrease strictly when
traced backwards, eventually giving a contradiction. In particular,
deleting vertices cannot circumvent the obstruction.

\subsubsection{1. A small-target growth
lemma}\label{a-small-target-growth-lemma}

\textbf{Lemma 1.} Let \(Q\) be a bipartite graph with parts \(X,Y\) and
girth greater than \(2M\). Suppose that nonempty sets \(P\subseteq X\)
and \(R\subseteq Y\) satisfy \[
|R|\le M
\] and every vertex of \(P\) has at least two neighbors in \(R\). Then
\[
|P|\le |R|-1.
\]

\textbf{Proof.} Any cycle in \(Q[P,R]\) would have length at most
\(2|R|\le 2M\), because a bipartite cycle uses equally many distinct
vertices from each part. Consequently, \(Q[P,R]\) is a forest. Therefore
\[
2|P|
\le e(Q[P,R])
\le |P|+|R|-1,
\] which proves the claim. \(\square\)

We will use regular bipartite graphs of arbitrarily large prescribed
girth. A fully explicit, self-contained construction is given in Section
4.

\subsubsection{2. Construction of the Eulerian
host}\label{construction-of-the-eulerian-host}

Fix \(d\ge 2\), and put \[
M=\binom d2,\qquad L=M+1.
\] Choose a simple \(d\)-regular bipartite graph \(Q\), with \(n\)
vertices in each part, whose girth is greater than \(2M\).

Take disjoint vertex sets \[
B,A_1,A_2,\ldots,A_L,
\qquad |B|=d,\qquad |A_i|=n.
\] Define \(D_d\) by the following arcs:

\begin{enumerate}
\def\labelenumi{\arabic{enumi}.}
\tightlist
\item
  All arcs from \(B\) to \(A_1\).
\item
  For each \(1\le i<L\), a copy of \(Q\) between \(A_i\) and
  \(A_{i+1}\), with every edge directed from \(A_i\) to \(A_{i+1}\).
\item
  All arcs from \(A_L\) to \(B\).
\end{enumerate}

There are no other arcs. Schematically, \[
B\ \longrightarrow\ A_1\ \overset{Q}{\longrightarrow}\ A_2\
\overset{Q}{\longrightarrow}\ \cdots\
\overset{Q}{\longrightarrow}\ A_L\ \longrightarrow\ B.
\]

Every vertex of \(B\) has in-degree and out-degree \(n\). Every vertex
of every \(A_i\) has in-degree and out-degree \(d\). Since \(n\ge d\),
\[
d^-_{D_d}(z)=d^+_{D_d}(z)\quad\text{for all }z,
\qquad
\delta^+(D_d)=d.
\] Also, \(L\ge2\), so this is an oriented graph: there are no
antiparallel pairs.

The digraph is strongly connected. Every vertex in an \(A_i\) can reach
\(B\) by following the layers forwards. Conversely, tracing predecessors
backwards through the regular layers shows that every vertex in an
\(A_i\) is reachable from \(A_1\), and hence from every vertex of \(B\).
Thus the construction satisfies even the convention that an Eulerian
digraph must be connected.

\subsubsection{\texorpdfstring{3. No \(F\)-free subdigraph can have
minimum out-degree
\(2\)}{3. No F-free subdigraph can have minimum out-degree 2}}\label{no-f-free-subdigraph-can-have-minimum-out-degree-2}

Suppose, for a contradiction, that a nonempty subdigraph
\(H\subseteq D_d\) is \(F\)-free and satisfies \[
\delta^+(H)\ge2.
\] Set \[
T=V(H)\cap B,\qquad S_i=V(H)\cap A_i.
\]

All these sets are nonempty. Indeed, every arc goes to the next class in
the cyclic sequence \[
B,A_1,\ldots,A_L,B.
\] Starting at any vertex of \(H\) and repeatedly following an outgoing
arc therefore visits every class.

\paragraph{The final layer is small}\label{the-final-layer-is-small}

For each \(z\in S_L\), choose two distinct out-neighbors of \(z\) in
\(H\). Both belong to \(T\); denote their unordered pair by \(P_z\).

If two distinct vertices \(z,z'\in S_L\) had \(P_z=P_{z'}\), these two
sources and their two common heads would form a copy of
\(\vec K_{2,2}\). Hence the pairs \(P_z\) are all distinct, and \[
|S_L|
\le \binom{|T|}{2}
\le \binom d2
=M.
\qquad\text{(3.1)}
\]

This is the elementary pair-counting obstruction from the supplied lead,
now applied at a designated final layer.

\paragraph{Every preceding layer is strictly
smaller}\label{every-preceding-layer-is-strictly-smaller}

Suppose \(1\le i<L\) and \(|S_{i+1}|\le M\). Every vertex of \(S_i\) has
at least two out-neighbors in \(S_{i+1}\), since these are its only
possible out-neighbors in \(H\).

The underlying graph between \(A_i\) and \(A_{i+1}\) is a copy of \(Q\).
Lemma 1 therefore gives \[
|S_i|\le |S_{i+1}|-1.
\qquad\text{(3.2)}
\]

Starting from (3.1) and applying (3.2) backwards through the \(L-1=M\)
intermediate transitions yields \[
|S_1|
\le |S_L|-(L-1)
\le M-M
=0,
\] contradicting the nonemptiness of \(S_1\).

This proves the theorem. Notice that \(H\) was not assumed to be
spanning or Eulerian.

If an Eulerian-avoidability function \(d_F\) existed, choosing \[
d\ge\max\{2,d_F(2)\}
\] and applying its defining property to \(D_d\) would contradict the
theorem. Therefore \(\vec K_{2,2}\) is not Eulerian-avoidable.
\(\square\)

\subsubsection{4. An explicit construction of the high-girth
layers}\label{an-explicit-construction-of-the-high-girth-layers}

For completeness, the required regular bipartite graph can be
constructed without invoking any external existence theorem.

\textbf{Lemma 2.} For every \(d\ge2\) and \(g\ge2\), one can explicitly
construct a finite simple \(d\)-regular bipartite graph of girth greater
than \(g\).

\textbf{Proof.} Let \(W\) consist of the empty word and all words of
length at most \(g\) over the alphabet \([d]\) in which consecutive
letters differ.

For each \(j\in[d]\), define a permutation \(\sigma_j\) of \(W\) as
follows:

\begin{itemize}
\tightlist
\item
  If a word ends in \(j\), delete its final letter.
\item
  Otherwise, if its length is less than \(g\), append \(j\).
\item
  Otherwise, fix it.
\end{itemize}

Each \(\sigma_j\) is an involution. We use right actions for permutation
multiplication.

If \[
j_1j_2\cdots j_\ell,\qquad 1\le\ell\le g,
\] has consecutive letters distinct, then \[
\varnothing\,\sigma_{j_1}\sigma_{j_2}\cdots\sigma_{j_\ell}
=j_1j_2\cdots j_\ell.
\] In particular, \[
\sigma_{j_1}\sigma_{j_2}\cdots\sigma_{j_\ell}\ne 1.
\qquad\text{(4.1)}
\]

Let \(\Gamma=\operatorname{Sym}(W)\). Form a bipartite graph \(Q\) whose
two parts are copies \(\Gamma_0,\Gamma_1\) of \(\Gamma\), with edges \[
\pi_0\,(\pi\sigma_j)_1
\qquad
(\pi\in\Gamma,\ j\in[d]).
\] The generators \(\sigma_1,\ldots,\sigma_d\) are distinct, so \(Q\) is
simple and \(d\)-regular. Give each displayed edge color \(j\). Each
color class is a perfect matching.

Suppose \(Q\) had a simple cycle of length \(\ell\le g\). Reading its
edge colors in order gives a sequence \(j_1,\ldots,j_\ell\) with
consecutive colors distinct: using the same color twice consecutively
would immediately traverse the same matching edge backwards. Since every
\(\sigma_j\) is an involution, returning to the starting vertex gives \[
\sigma_{j_1}\sigma_{j_2}\cdots\sigma_{j_\ell}=1,
\] contrary to (4.1). Hence \(Q\) has girth greater than \(g\).
\(\square\)

For the counterexample, take \(g=2M=d(d-1)\). This gives a completely
specified family. More explicitly, one may use \[
|W|=1+d\sum_{r=0}^{d(d-1)-1}(d-1)^r,
\qquad n=|W|!,
\] so that \[
|V(D_d)|=d+\left(\binom d2+1\right)n.
\] These sizes are deliberately unoptimized.

\subsubsection{5. Scope and a
generalization}\label{scope-and-a-generalization}

The same construction proves a broader statement:

\textbf{Corollary.} For all \(s,t\ge2\), the one-directed complete
bipartite digraph \(\vec K_{s,t}\), directed from its \(s\)-vertex part
to its \(t\)-vertex part, is not Eulerian-avoidable.

To see this, fix \(d\ge t\), use target minimum out-degree \(k=t\), and
replace the parameters by \[
M=(s-1)\binom dt,\qquad L=M+1.
\] In a \(\vec K_{s,t}\)-free subdigraph, choosing a \(t\)-set of
out-neighbors for each final-layer vertex assigns any particular
\(t\)-set to at most \(s-1\) sources. Thus the final layer again has at
most \(M\) vertices. Since \(t\ge2\), Lemma 1 again forces strict growth
forwards through every small intermediate layer, giving the same
contradiction.

For the original question, only \(s=t=2\) is needed. The crucial
distinction from complete cyclic blow-ups is that the intermediate
high-girth layers cannot support a small \(2\)-out core. This removes
the vertex-deletion escape available in the supplied attempt and gives a
counterexample to Eulerian-avoidability itself, not merely to a spanning
version.

\appendix

\section{Provenance and review
status}\label{provenance-and-review-status}

\begin{warning}

\textbf{UNREFEREED MODEL OUTPUT.} This document was selected solely
because GPT-6 Astra labelled its own result
\texttt{would\_publish:\ true} and returned \texttt{proved} or
\texttt{disproved}. It has not passed the adversarial LLM referee used
for the earlier Sol campaign, has not been checked by a human
mathematician, and has not been checked for novelty. Treat every
mathematical and bibliographic claim as unverified.

\end{warning}

{\def\LTcaptype{none} % do not increment counter
\begin{longtable}[]{@{}
  >{\raggedright\arraybackslash}p{(\linewidth - 2\tabcolsep) * \real{0.5000}}
  >{\raggedright\arraybackslash}p{(\linewidth - 2\tabcolsep) * \real{0.5000}}@{}}
\toprule\noalign{}
\endhead
\bottomrule\noalign{}
\endlastfoot
Catalog id & \texttt{2510.11311\_\_04} \\
Catalog entry &
\href{https://graph-theory-ai.github.io/graph-conjectures/arxiv/2510.11311__04/}{Question
6.4} \\
Source corpus & arxiv \\
Campaign leg & selected second attempts (\texttt{attacks\_retry}) \\
Source paper / entry & Extending Thomassen's conjecture to directed
graphs \\
Model verdict & \textbf{disproved} (confidence: high) \\
Selection criterion & \texttt{would\_publish:\ true} and verdict
\texttt{proved} or \texttt{disproved} \\
Model & \texttt{gpt-6-astra}, reasoning effort \texttt{max},
\texttt{mode=pro}, flex service tier \\
Original artifact &
\texttt{attacks\_retry/2510.11311\_\_04/output.md} \\
Independent review & \textbf{None} \\
arXiv & \href{https://arxiv.org/abs/2510.11311}{arXiv:2510.11311} \\
Previous attempt & \texttt{attacks/2510.11311\_\_04}: gpt-5.6-sol
verdict \texttt{partial} \\
Model's caveats & The counterexamples have very large degree spread;
this does not classify the other orientations of C4. \\
\end{longtable}
}

The generated Markdown and LaTeX sources for this PDF are stored under
\texttt{to\_review\_astra/src/2510.11311\_\_04/}.

\end{document}
